Our analysis pipeline is illustrated in Figure \ref{fig:workflow_two_column_scaled}. We take the following steps.

\begin{figure}[htbp]     % full-page figure
\centering
\begin{adjustbox}{center,max width=\textwidth}
\begin{tikzpicture}[
    node distance = 1.6cm and 0.8cm,
    every node/.style = {
        rectangle,
        draw,
        rounded corners,
        align=center,
        text width=5cm,
        minimum height=1.3cm,
        font=\small
    },
    >={Stealth[length=6pt,width=8pt]}
]

% -----------------------------------------------------
% Stage 1 (left column)
% -----------------------------------------------------
\node (bench)  {Benchmark sample\\of human-labeled transcripts\\\small(\textit{N = \data{summary_stats/tagging_performance/benchmark_sample_count_int}})};
\node (dev)    [below=of bench] {Develop \& compare\\prompts \& LLM models};
\node (best)   [below=of dev] {Select best-performing\\prompt + model};

% -----------------------------------------------------
% Stage 2 (right column)
% -----------------------------------------------------
\node (full)   [right=5.6cm of bench] {Apply best model to\\Full transcript dataset\\(Capital IQ)\\\small(\textit{N = \data{summary_stats/dataset_overview/total_transcripts_int}})};
\node (flag)   [below=of full] {LLM-flagged transcripts\\score $\ge 75$\\\small(\textit{N = \data{summary_stats/tagging_performance/llm_tagged_count_int}})};
\node (valid)  [below=of flag] {LLM-validated transcripts\\11-query avg $\ge 75$\\\small(\textit{N = \data{summary_stats/tagging_performance/llm_validation_tagged_int}})};
\node (audit)  [below=of valid] {Human-audited\\LLM-validated transcripts\\\small(\textit{N = \data{summary_stats/tagging_performance/human_audit_sample_count_int}})};

% -----------------------------------------------------
% Arrows within columns and between stages
% -----------------------------------------------------
\draw[->, thick] (bench) -- (dev);
\draw[->, thick] (dev)   -- (best);

\draw[->, thick] (full)  -- (flag);
\draw[->, thick] (flag)  -- (valid);
\draw[->, thick] (valid) -- (audit);

\draw[->, thick] (best.east) -- ++(0.8cm,0) |- (full.west);

% -----------------------------------------------------
% Stage boxes (in background with extra padding)
% -----------------------------------------------------
\begin{scope}[on background layer]
    \node[draw,dashed,rounded corners,
          inner ysep=0.55cm, inner xsep=0.4cm,
          fit=(bench)(best)] (s1box) {};
    
    \node[draw,dashed,rounded corners,
          inner ysep=0.55cm, inner xsep=0.4cm,
          fit=(full)(audit)] (s2box) {};
\end{scope}

% -----------------------------------------------------
% Stage titles
% -----------------------------------------------------
\node[font=\bfseries] at (s1box.north) [yshift=1.0cm] {Stage 1: Prompt \& Model Development};
\node[font=\bfseries] at (s2box.north) [yshift=1.0cm] {Stage 2: Large-Scale Analysis};

\end{tikzpicture}
\end{adjustbox}
\caption{Analysis Workflow}
\label{fig:workflow_two_column_scaled}
\begin{minipage}{\textwidth}
    \vspace{1em}
    \footnotesize
    \textit{Notes:} This figure illustrates the workflow for detecting collusive corporate communications. Stage 1 develops and evaluates prompts and models using the human-labeled benchmark dataset through cross-validation to select the best-performing combination. Stage 2 applies the chosen model at scale to the full Capital IQ transcript dataset (\data{summary_stats/dataset_overview/total_transcripts_int} transcripts), refines high-scoring transcripts with multiple queries using a 75-point threshold, and conducts human audits of the top-scoring results to validate model performance.
\end{minipage}
\end{figure}


\textbf{Construction of human-labeled benchmark sample}

We constructed a human-labeled benchmark sample using two sources. 
The first is the \textcite{aryal2022coordinated} transcripts, with 
their binary assessment of collusion. The second are transcripts 
from U-Haul's parent company, Amerco. Some of these transcripts 
were previously discussed and considered collusive by \textcite{harrington2022collusion}. 
Citing these communications, the FTC charged U-Haul with an ``invitation to collude'' 
in \textcite{uhaul2010ftc}, resulting in a settlement. We reviewed 
these transcripts ourselves and provided collusion severity scores 
on a 0-100 scale. To obtain a binary classification for the U-Haul 
sample, we set a threshold value of 50. This leaves us with a 
benchmark sample of \data{summary_stats/tagging_performance/benchmark_sample_count_int} transcripts.

\textbf{Development and evaluation of prompts and LLM models}


With the benchmark sample in hand, we developed and evaluated 
prompts and experimented with different LLM models (see Appendix~\ref{sec:app_prompt_comparison} for prompt comparison).\footnote{Recent work by \textcite{asirvatham2026gpt} supports our approach in prompting, variable extraction, and benchmarking through a more generalized package that uses OpenAI LLMs on a selection of tasks. They also highlight using such LLMs as ``measurement tools'' and showcase ways to improve reliability and validity of LLM output.} We ultimately selected \textit{gpt-4o-mini} (temperature=1.0) and the following zero-shot prompt:
\begin{quote}
    \textit{You’re an expert in detecting subtle collusive intent in company earnings calls. You know firms can try to avoid scrutiny if they express intentions to reduce supply or increase prices, or they can imply it would be good if industry capacity were limited or prices were kept high. Analyze the following transcript of a public company communication and identify any signs that the company is trying to coordinate with competitors to limit capacity or keep prices high. Provide a score from 0 to 100, where 0 means no evidence of collusive intent and 100 means strong evidence of collusive intent. Additionally, provide a very brief explanation of your assessment. Also provide a list of excerpts that support your assessment.}
\end{quote}

This prompt was crafted to mention the broad types of collusive actions 
(i.e. maintaining high prices, limiting capacity), and specify 
clearly that we target a company's intention to coordinate on such actions.

\textbf{Large-scale analysis of Capital IQ transcripts}


We applied the best-performing prompt and model to the full 
Capital IQ transcript dataset (N=\data{summary_stats/dataset_overview/total_transcripts_int}). 
We define a transcript as {\textbf{LLM-flagged}} if it has an initial 
collusion score of at least 75. These are transcripts where the initial 
LLM assessment was that there is significant evidence of collusive intent. 
This flagged \data{summary_stats/tagging_performance/llm_tagged_count_int} transcripts.

\textbf{LLM Validation of LLM-flagged transcripts}


Because the LLM we use generates non-deterministic outputs for the same 
given transcript/prompt pair, a single evaluation of a transcript may 
reflect idiosyncratic sampling noise rather than a stable assessment. 
To address this issue and to produce a more robust output, we ran an 
additional ten repeated queries on the LLM-flagged transcripts. We then 
define as {\textbf{LLM-validated}} transcripts those whose average 
score remained at least 75 after the validation. This leaves us with
 \data{summary_stats/tagging_performance/llm_validation_tagged_int} transcripts.

\textbf{Human audit of LLM-validated transcripts}


Finally, we manually audited the top \data{summary_stats/tagging_performance/human_audit_sample_count_int} 
LLM-validated transcripts. The audit is crucial for both of the study's objectives: 
understanding the content used to facilitate coordinated conduct to reduce 
competition and evaluating the performance of an LLM as a screen for enforcers.
 The audit procedure is described in section \ref{sec:audit}.
